\documentclass{article}
%% ------------------------
%%
%% compile with:
%%  latex rodipos
%%  dvips -t poster rodipos 
%%
%% ``poster'' is a new media size which has to be defined in config.ps
%% of dvips. Information conform with the DSC about the required media
%% size will be included in the resulting .ps file. This allows tools
%%  operating on the .ps file like ghostview to extract the media size.
%%
%% config.ps is located in /etc/texmf/dvips 
%% We currently use the follwing definition:
%%
%%        @ poster 36in 100cm
%%        @+ ! %%DocumentMedia: Poster 2592 2834
%%        @+ ! %%DocumentPaperSizes: Poster
%%        @+ %%BeginPaperSize: Poster
%%        @+ %%EndPaperSize
%XW%
%%
%% How to distribute the poster on multiple A4 pages
%% -------------------------------------------------
%%
%% compile with:
%%  latex rodipos
%%  dvips -E rodipos
%%  poster -v -mA4 -pA0 rodipos.ps > rodipos-A4.ps
%%
%%  dvips generates an Encapsulated PostScript file. The dimensions of the
%%  EPS graphics depend on the .tex source! Note, that we do not specify
%%  a paper size in this case.
%%  ghhostview is able to display the resulting EPS file.
%%  
%%  In the next step we use the poster command to scale the EPS file to
%%  the desired final poster size (here A0) and distribute the output over
%%  several smaller pages (here A4). The poster command has a large number
%%  of options to specify the size of the desired output and the media
%%  size.
%%
%% References:
%% -----------
%% [1] Beulshausen, Guenther (1998) Druck-TeXnik. iX 9/1998 pp 130--131
%% [2] van Eijndhoven, Jos (1995) poster man page, available on CTAN
%%
%%


%%%%%%%%%%%%%%%%%%%%%%%%%%%%%%%%%%%%%%%%%%%%%%%%%%%%%%%%%%%%%%%%%%%%%%%%
%% Arbitrary scaling of TeX fonts
%% ------------------------------
%%
%% Package type1cm, allows TeX fonts to be scaled to (almost) arbitrary
%% sizes. Use this package with care, because for fonts which are not
%% scalabe (like PostScript fonts), a large amount of font data will be
%% generated.
%%
%% Support for ams is available in type1cm, however
%% amsmath needs to be included before type1cm
%%
\usepackage{mathtools}
\usepackage{palatino}
\usepackage{amssymb}
\usepackage{type1cm}
\usepackage{pstricks}
\usepackage{pst-node}
\usepackage[skip=2pt,font=Large]{caption}

\usepackage{float}
\usepackage{graphicx}
\usepackage{array}
\usepackage{url}
\usepackage{enumitem}

\usepackage{tikz}
\usepackage{pgfplots}
\pgfplotsset{compat=1.11}
\usepgfplotslibrary{external} 
\tikzexternalize
%\tikzset{external/force remake}
\usepackage{cite}
\special{papersize=36in,120cm}
%\special{papersize=53in,45in}

\renewcommand{\familydefault}{cmss}
\renewcommand{\figurename}{Fig.}

%%%%%%%%%%%%%%%%%%%%%%%%%%%%%%%%%%%%%%%%%%%%%%%%%%%%%%%%%%%%%%%%%%%%%%%%
%% Arbitrary page size
%% -------------------
%%
%% The package geometry makes the definition of arbitrary pages sizes
%% more conveniant.
%%
%% Delmenhorst00 maximal poster is width 120cm, height 140cm
%%
%% HP DJ2500 paper width is 36in
%%
%%    36in * 2.54cm/in = 91.44cm
%%    
%% ==> 91.44cm x 100cm fits
%%
\usepackage{calc}
%\usepackage[paperwidth=36in,paperheight=120cm,margin=1in]{geometry}
\usepackage[paperwidth=32in,paperheight=120cm,margin=1in]{geometry}
%\usepackage[paperwidth=53in,paperheight=45in,textwidth=1175mm,textheight=36in]{geometry}
%\usepackage{lscape}

%%%%%%%%%%%%%%%%%%%%%%%%%%%%%%%%%%%%%%%%%%%%%%%%%%%%%%%%%%%%%%%%%%%%%%%%
%% Further packages and adjustments
%% --------------------------------
%%
%%
%% printer HP DeskJet 2500 left margin is 1in --> \hoffset = -1in

\usepackage{multicol}
\setlength{\multicolsep}{0pt}
\setlength{\columnsep}{0.75in}
\setlength{\columnseprule}{1pt}
\setlength{\parindent}{0cm}

\pagestyle{empty}

\newcommand{\VisibleStrut}[2][0cm]{\rule[#1]{1mm}{#2-#1}}
\newcommand{\Strut}[2][0cm]{\rule[#1]{0mm}{#2-#1}}


\newpsobject{showgrid}{psgrid}{subgriddiv=2,griddots=10,gridlabels=6pt}

%%%%%%%%%%%%%%%%%%%%%%%%%%%%%%%%%%%%%%%%%%%%%%%%%%%%%%%%%%%%%%%%%
%%
%% private macros
%%
\renewcommand{\sin}{\mathsf{sin}}
\renewcommand{\arctan}{\mathsf{arctan}}
\renewcommand{\cos}{\mathsf{cos}}
\renewcommand{\exp}{\mathsf{exp}}
\renewcommand{\log}{\mathsf{log}}
\renewcommand{\max}{\mathsf{max}}
\renewcommand{\sup}{\mathsf{sup}}

\newcommand{\note}[1]{\mbox{}\marginpar{\raggedleft\hspace{0pt}#1}}
\newcommand{\inquotes}[1]{``#1''}
\newcommand{\synod}{\textsf{SYNOD}}
\newcommand{\ppSet}{\mathsf{N}}
\newcommand{\Set}[1]{\mathsf{#1}}
\newcommand{\pp}{\mathsf{a}} % point process
\newcommand{\PP}{\mathsf{A}} % point process
\newcommand{\cm}{n} % counting measure
\newcommand{\ms}{\text{ ms}} % millisecond unit in math mode
\newcommand{\diff}{\operatorname{d}}
\newcommand{\T}{\mathsf{T}}
\newcommand{\R}{\mathbb{R}_0^+}
\newcommand{\N}{\mathbb{N}}
\newcommand{\OP}{\boldsymbol{\Phi}}
\newcommand{\indicator}{\boldsymbol{1}}
\newcommand{\Hist}{\mathcal{H}}
\newcommand{\E}[1]{\bigl<#1\bigr>}
\newcommand{\erf}{\operatorname{erf}}
\newcommand{\var}{\operatorname{var}}
\newcommand{\SD}{\operatorname{SD}}
\newcommand{\op}{\circ}
\newcommand{\sfc}{\mathsf{S}}
\newcommand{\psp}{\operatorname{psp}}
\newcommand{\const}{\operatorname{const}}
\newcommand{\MRef}{\mathcal{R}}

\newcommand{\Secref}[1]{Sec.~\ref{#1}}

\newcommand{\chapref}[1]{\ref{#1}}
\newcommand{\Chapref}[1]{Chap.~\ref{#1}}
\newcommand{\cpanel}[1]{{\color{blue}\textbf{#1} }}



%%%%%%%%%%%%%%%%%%%%%%%%%%%%%%%%%%%%%%%%%%%%%%%%%%%%%%%%
%% Macros for cross-references, and citations
%%
%%

\newcommand{\figref}[1]{\ref{#1}}
\newcommand{\Figref}[1]{\textcolor{blue}{Fig.~#1}}

\newcommand{\figpartref}[2]{\ref{#1}~\textbf{#2}}
\newcommand{\Figpartref}[2]{Fig.~\ref{#1}~\textbf{#2}}

\newcommand{\panel}[1]{\textbf{#1}}

\newcommand{\parenthesizedfrac}[2]{\genfrac{(}{)}{}{}{#1}{#2}}

\newcommand{\FigLabel}[0]{%
{\raggedright%
%\colorbox{gray}{\color{blue}\Figref{PacketPotentialApproachFigure}}
\psframebox[fillstyle=solid,fillcolor=gray,linecolor=gray]
{\stepcounter{figure}\textcolor{blue}{Figure~\arabic{figure}}\Strut[-1cm]{0cm}}}
\vspace{3mm}}

\definecolor{salmon}{rgb}{1,0.6,0.6}

%%%%%%%%%%%%%%%%%%%%%%%%%%%%%%%%%%%%%%%%%%%%%%%%%%%%%%%%
%% Macros for Operator and other special symbols
%%
%%
%%
\newcommand{\eqfullstop}{\quad\text{.}}
\newcommand{\eqcomma}{\quad\text{,}}

%
% figure style
%

% tikz figure settings
\newlength\figureheight
\newlength\figurewidth


\newsavebox{\fb}
\newlength{\wf}
\newlength{\hf}

\newlength{\FigureLabelMargin}
\newlength{\FigureXLabelDistance}
\setlength{\FigureLabelMargin}{-0.5cm}
\setlength{\FigureXLabelDistance}{-0.05cm}

\newcommand{\LabelWithUnit}[2]{#1\quad(\! #2)}
\newcommand{\PanelLabel}[1]{\fontsize{8mm}{8mm}\bfseries #1}


\newcommand{\lbox}[1]{%
{\protect\footnotesize%
\protect\setlength{\protect\fboxsep}{0.3ex}%
\protect\raisebox{0.2ex}{%
\protect\fbox{\protect#1}}%
}%
}


% Gleichungen
\newcommand{\beq}{\begin{equation}}
\newcommand{\eeq}{\end{equation}}
\newcommand{\beqo}{\begin{displaymath}}
\newcommand{\eeqo}{\end{displaymath}}
\newcommand{\bea}{\begin{eqnarray}}
\newcommand{\eea}{\end{eqnarray}} 
\newcommand{\beao}{\begin{eqnarray*}}
\newcommand{\eeao}{\end{eqnarray*}}
% Doppelte und einfache Anführungszeichen
\newcommand{\afz}[1]{``#1''}
\newcommand{\afzsimple}[1]{`#1'}

\makeatletter

\renewcommand{\section}{%
\@startsection{section}{1}{0pt}{3ex}{1ex}{%
\color[rgb]{0,0.4078,0.7059}%\color{blue}
\fontsize{13mm}{13mm}\sffamily\bfseries}}

\renewcommand{\subsection}{%
\@startsection{subsection}{1}{0pt}{2ex}{1ex}{%
\color[rgb]{0,0.4078,0.7059}
\fontsize{11mm}{11mm}\sffamily\bfseries}}

\makeatother 

\begin{document}

\bibliographystyle{plain}
\rule{\linewidth}{1pt}
\vspace*{1in}

% ======================================================
% poster head
% ======================================================
\begin{center}
{%
\fontsize{2.4cm}{2.4cm}\sffamily
Interactive Online Learning \\for Obstacle
Classification on a Mobile Robot
\\[5mm]
\fontsize{1.3cm}{1.3cm}\sffamily%
Viktor Losing$^{1,2}$, Barbara Hammer$^1$, Heiko Wersing$^{2}$ \\[5mm]
\fontsize{1.3cm}{1.3cm}\sffamily%
$1$ -- Research Institute for Cognition and Robotics - CoR-Lab at Bielefeld University,\\  $2$ -- HONDA Research Institute Europe GmbH\\
\vspace*{1in}
\rule{\linewidth}{1pt}
\fontsize{8mm}{9mm}\sffamily}
\end{center}


\fontsize{10mm}{11.25mm}\sffamily
\begin{multicols}{3}

\include{newcommands}

\addtolength{\abovedisplayskip}{10pt}
\addtolength{\belowdisplayskip}{10pt}

\newcommand{\fv}{{\mathbf f}}
\newcommand{\cv}{{\mathbf c}}

\parskip 1.0pc


% ======================================================
% Introduction
% ======================================================
\section*{Online Learning}
\begin{itemize}
	\item Motivation:
	\begin{itemize}
		\item Adaptation to user habits/environment
		\item Current methods rather simple
	\end{itemize}	
	\item Benefits:
	\begin{itemize}
		\item Incremental learning
		\item Incorporation of new classes
		\item Control of model complexity
		\item Handling of drift
	\end{itemize}
	\item Challenges:
	\begin{itemize}
		\item Stability - Plasticity
	\end{itemize}		
\end{itemize}

\section*{Interactive Scenario}
\newcommand{\myCaption}{\em \fontsize{8mm}{9.25mm}\sffamily}

%\renewcommand*\labelitemi{\normalfont\bfseries\textendash}
\begin{itemize}
	\item Random exploration of garden environment
	\item Grass-segmentation for obstacle detection
	\item Live-labeling via iPad
	\item Object specific actions are performed
	\item Confidence estimation to identify unknown objects \cite{CN09} 
\end{itemize}

\begin{figure}[H]
\begin{center}
\includegraphics[width=0.3\textwidth]{Images/General/Scenario.jpg}
\end{center}
\begin{center}
{\myCaption Fig.~1.~Typical scene of the interactive scenario.}
\end{center}
\end{figure}


\begin{figure}[H]
\begin{center}
\includegraphics[width=0.1635\textwidth]{Images/General/VisualField.png}
\hspace{10 pt}
\includegraphics[width=0.135\textwidth]{Images/General/TallObject.jpg}
\end{center}
{\myCaption Fig.~2.~Left: Partition of the input image. Objects are only examined as long
as they are in the region of interest. Overlaps with the vision bump region end
the examination and trigger a reaction. Information within the top marked area
is ignored completely. Right: Masked area (red) of a close tree stump. The
upper part of the image is disregarded even though it belongs to the object.}
\end{figure}


\subsection*{Feature Representation}
\begin{minipage}[c]{0.57\linewidth}
\begin{itemize}
	\item Color-based (simple$+$robust)
	\item Rg-Chromaticity histogram
	\item Intensity invariant
	\item Size invariant
	\item 21 dimensions	
\end{itemize}

\end{minipage}
\begin{minipage}[c]{0.4\linewidth}
\begin{figure}[H]
\includegraphics[width=1\linewidth]{Images/General/RgChromacityQuads.png}		
{\myCaption Fig.~3.~Histogram binning of the rg-chromaticity space.}
\end{figure}
\end{minipage}



\begin{figure}[H]
\begin{center}
\includegraphics[width=0.15\textwidth]{Images/General/BlueBall_4_10.jpg}
\hspace{10 pt}
\includegraphics[width=0.15\textwidth]{Images/General/BlueBall_4_10_rgComplete_bright.jpg}
\end{center}
\begin{center}
{\myCaption Fig.~4.~Image in RGB and rg-chromaticity space.}
\end{center}
\end{figure}



\section*{Outdoor Benchmark Data-set}
\begin{minipage}[c]{0.6\linewidth}
\begin{itemize}
	\item 40 objects
	\item 10 sequences per object
	\item 10 images per sequence
	\item Cloudy \& sunny conditions
\end{itemize}

\end{minipage}
\begin{minipage}[c]{0.4\linewidth}
\begin{figure}[H]
\begin{center}
\includegraphics[width=1\linewidth]{Images/General/Runs.jpg}
\end{center}
{\myCaption Fig.~5.~First, fifth and tenth image of some sequences.}

\end{figure}

\end{minipage}

\begin{figure}[H]
\begin{center}
\includegraphics[width=0.31\textwidth]{Images/General/Challenges.png} 
\end{center}
{\myCaption Fig.~6.~Challenges of the Outdoor data set. Objects covered in various degrees
by shadows cast by environmental obstacles such as trees or buildings (a)
or by the robot itself (b). Sunlight and auto-exposure cause a darker object
representation (c). Specular highlights (d). Occlusions by the image border
(e). Various object poses (f).}
\end{figure}



% ======================================================
% Methods
% ======================================================

\section*{Learning Architecture}

\begin{minipage}[c]{0.6\linewidth}
\begin{itemize}
	\item Generalized LVQ \cite{SO95} 
	\item Supervised, prototype based
	\item $E(X, W)=\sum_{i=1}^{m} \Phi(\frac{(d^+_i)-d^-_i))}{(d^+_i+d^-_i))})$
	\item Minimized in stochastic gradient descent scheme
\end{itemize}

\end{minipage}
\begin{minipage}[c]{0.4\linewidth}
\begin{figure}[H]
\begin{center}
\includegraphics[width=1\linewidth]{Images/General/learning21.png}
\end{center}
{\myCaption Fig.~7.~Simple three class problem. Black lines represent the class borders.}
\end{figure}

\end{minipage}

\subsection*{Prototype Insertion - Current Methods}


\begin{minipage}[c]{0.58\linewidth}
\begin{itemize}
	\item Based on heuristics
	\item Use misclassifications only
\end{itemize}

\end{minipage}
\begin{minipage}[c]{0.4\linewidth}
\begin{figure}[H]
\begin{center}
\includegraphics[width=1\linewidth]{Images/General/Insertions1.png}
\end{center}
{\myCaption Fig.~8.~Situation before a new prototype is inserted.}
\end{figure}
\end{minipage}

\begin{figure}[H]
\includegraphics[width=1\linewidth]{Images/General/PlacementHeuristics.png}
\begin{center}
{\myCaption Fig.~9.~Situation after a new prototype is inserted by each heuristic.} 
\end{center}
\end{figure}


\subsection*{Prototype Insertion - SamplingCost}
\begin{itemize}
	\item Store a window of recent t samples in a list $\Psi$
	\item Choose a random subset $\hat{\Psi}\subseteq \Psi$ of size $|\hat{\Psi}|=\hat{t}$
	\item Evaluate each sample $(x_i,y_i)$ in $\hat{\Psi}$ as prototype:
	\begin{itemize}
	\item Add sample as prototype $\hat{W_i}:= \{W\cup (x_i,y_i)\ \}$
	\item Calculate the cost $E(\Psi, \hat{W_i})$
	\end{itemize}
	\item Choose $\hat{W_i}$ s.t. $E(\Psi, \hat{W_i})$ is minimized
\end{itemize}

\section*{Results}
\subsection*{Artificial Data}
\begin{figure}[H]
\begin{center}
\includegraphics[width=0.31\textwidth]{Images/Toy/BorderDS.png}
\end{center}
\begin{center}
\includegraphics[width=0.31\textwidth]{Images/Toy/OverlapDS.png}
\end{center}
{\myCaption Fig.~10.~Networks of each placement strategy after the training of the data
sets Border (top) and Overlap (bottom).}
\end{figure}

\begin{minipage}[c]{0.48\linewidth}

\centering
\fontsize{8.5mm}{9.25mm}\sffamily
\begin{tabular}{c|cc}
\textit{Border-DS} & Acc. & Nodes\\ \hline \\ [-1.5ex]
SamplingCost & \textbf{93.51} & \textbf{38.2}\\
Closest & 90.17 & 58.2\\
Cluster & 91.93 & 42.9\\
Voronoi & 91.71 & 46.4
\end{tabular}
\end{minipage}\hspace{20 pt}
\begin{minipage}[c]{0.48\linewidth}
\centering
\fontsize{8.5mm}{9.25mm}\sffamily
\begin{tabular}{c|cc}
\textit{Overlap-DS} & Acc. & Nodes\\ \hline \\ [-1.5ex]
SamplingCost & \textbf{78.74} & \textbf{21.3}\\
Closest & 65.76 & 29.8\\
Cluster & 74.85 & 25.6\\
Voronoi & 74.08 & 26.4
\end{tabular}
\end{minipage}

\subsection*{Real Data}
\begin{figure}[H]
\includegraphics[width=0.31\textwidth]{Images/Outdoor/OutdoorPrototypeDistribution.png}
\begin{center}
{\myCaption Fig.~11.~Prototype distributions for the Outdoor data set.} 
\end{center}
\end{figure}

\begin{center}
\begin{tabular}{c|ccc}
\fontsize{8.5mm}{9.25mm}\sffamily
\textit{Random-order} & Test acc. & Train acc. & Nodes\\ \hline \\ [-1.5ex]
\rule{0pt}{10pt}
SamplingCost & \textbf{81.58} & 85.94 &\textbf{234.8}\\
Closest & 78.03 & 83.23 & 253.4\\
Cluster & 78.55 & 83.06 & 247.0\\
Voronoi & 81.18 & 85.62 & 236.6
\end{tabular}
\end{center}

\begin{center}
\fontsize{8.5mm}{9.25mm}\sffamily
\begin{tabular}{c|ccc}
\textit{Sequence-order} & Test acc. & Train acc. & Nodes\\ \hline \\ [-1.5ex]
\rule{0pt}{10pt}
SamplingCost & \textbf{61.38} & 91.04& 232.0\\
Closest & 59.58 & 88.04 & 245.8\\
Cluster & 59.18 & 88.96 & 235.6\\
Voronoi & 58.81 & 89.22 & \textbf{231.3}
\end{tabular}
\end{center}

\begin{figure}[H]
\begin{center}
\includegraphics[width=0.31\textwidth]{Images/General/DifficultObjects.png}
\end{center}
\begin{center}
{\myCaption Fig.~12.~Objects requiring the highest amount of prototypes.} 
\end{center}
\end{figure}

\begin{figure}[H]
\begin{center}
\includegraphics[width=0.31\textwidth]{Images/Outdoor/Easy.jpg}
\end{center}
\begin{center}
{\myCaption Fig.~13.~The easiest objects.} 
\end{center}
\end{figure}



\subsection*{Comparison with Incremental SVM \cite{Dl13}}

\begin{center}
\fontsize{8.5mm}{9.25mm}\sffamily
\begin{tabular}{c|ccc}
\textit{} Outdoor& \shortstack{Acc. / Nodes  \\500 samples} & \shortstack{Acc. / Nodes  \\1500 samples} & \shortstack{Acc. / Nodes  \\2800 samples}\\ \hline \\ [-1.5ex]
\rule{0pt}{10pt}
SC-GLVQ & 42.9 / 116.4 & 58.6 / 446.2& 64.3 / 948.8\\
SC-GMLVQ & \textbf{44.1 / 100.4} & \textbf{59.3 / 376.6} & \textbf{67.1 / 773.5}\\
iSVM & 40.6 / 363.3 & 57.5 / 777.0 & 65.2 / 1347\\
\end{tabular}
\end{center}

\begin{center}
\fontsize{8.5mm}{9.25mm}\sffamily
\begin{tabular}{c|ccc}
\textit{} COIL-100& \shortstack{Acc. / Nodes  \\500 samples}  & \shortstack{Acc. / Nodes  \\1000 samples}  & \shortstack{Acc. / Nodes  \\1700 samples}\\ \hline \\ [-1.5ex]
\rule{0pt}{10pt}
SC-GLVQ & 83.5 / 337.8 & 90.4 / 560.0  & 93.6 / 810.0\\
SC-GMLVQ & \textbf{85.0 / 330.6} & \textbf{91.9 / 539.5}& \textbf{94.7 / 760.0}\\
iSVM & 85.6 / 1834  & 92.7 / 2139 & \textbf{95.6 / 2501}\\
\end{tabular}
\end{center}


% ======================================================
% literature
% ======================================================
\begin{thebibliography}{}

\bibitem{CN09} Fischer L., Hammer B. and Wersing H. Rejection Strategies for Learning Vector Quantization. ESANN (2014)
\bibitem{SO95}Sato A. and Yamada, K. Generalized Learning Vector Quantization. NIPS (1995) 
\bibitem{KS05}Kirstein S., Wersing H. and K{\"o}rner E. Rapid Online Learning of Objects in a Biologically Motivated Recognition Architecture. 27th Pattern Recognition Symposium DAGM (2005)
\bibitem{GC09}Grbovic M. and Slobodan V. Learning Vector Quantization with adaptive prototype addition and removal. IJCNN (2009)
\bibitem{BO98}Bermejo S., Cabestany J. and Payeras{-}Capell{\`{a}} A new dynamic LVQ-based classifier and its application to handwritten character recognition. ESANN (1998)
\bibitem{Dl13}Diehl C. and Cauwenberghs G. SVM Incremental Learning, Adaptation and Optimization. IJCNN (2013)
\end{thebibliography}

\begin{figure}[H]
\begin{center}
\includegraphics[width=0.1\textwidth]{logos/hri_logo.jpg}
\hspace{220 pt}
\includegraphics[width=0.1\textwidth]{logos/rgb-cor-lab-logo1_text.jpg}
\end{center}
\end{figure}


\end{multicols}
\end{document}